\makeabstract
{
Don’t Be Rude, Daisy: Socially-Aware Robot Navigation
}
{
Madalina Voica (1), Professor Scott Alfeld, PhD (1)
}
{
\textsuperscript{1}Department of Computer Science, Amherst College, Amherst, MA, USA
}
{
This project investigates how Boston Dynamics’ quadruped robot Spot, nicknamed Daisy, can be equipped with socially-aware navigation skills in human-populated environments. The goal is to enable Daisy to move respectfully through dynamic spaces by modeling behaviors that prevent crowding, support group integration, and adhere to human proxemics.
}
{
A real-time multi-person detection and tracking pipeline using YOLOv7 and StrongSORT was implemented to allow Daisy to perceive and identify individuals across frames. Full 360-degree perception was achieved by streaming all five onboard cameras in parallel using multithreading. Depending on social context, Daisy uses either the Social Forces Model (for individual interactions) or the BOIDs model (for group movement). The Social Forces equation uses dispersion constants to control attraction to goals, repulsion from humans, and avoidance of obstacles. The BOIDs model includes weights for separation, alignment, and cohesion, guiding Daisy to move within a group as an active participant rather than trailing behind.
}
{
Daisy is able to detect and consistently track people across multiple viewpoints and estimate their relative velocities and positions. The use of both SFM and BOIDs allows flexible navigation strategies depending on social context. Initial field tests show stable multi-camera tracking, successful person-following behavior, and consistent avoidance of personal space violations.
}
{
This study demonstrates the feasibility of applying hybrid navigation models to enable context-aware, socially compliant movement in a mobile quadruped robot. These behaviors were driven by manually chosen parameters, but future work will focus on optimizing these parameters through reinforcement learning in a simulated environment to improve adaptability and realism in navigation behavior.
}

\newpage
